\clearpage

\section{GNN을 이용한 AST 기반의 취약점 탐지}

\subsection{개요}
Agent 모듈은 정적 분석 파이프라인의 후반부에서 동작하는 그래프 신경망(Graph Neural Network, GNN) 기반의 취약점 탐지 시스템이다. 이 모듈은 앞선 단계에서 생성된 코드의 그래프 표현(AST)을 입력으로 받아, 해당 함수가 취약한지(Vulnerable) 혹은 안전한지(Safe)를 분류하도록 학습된 PyTorch Geometric 모델을 운용한다.

\subsubsection{해결하고자 하는 문제}
기존의 정적 분석 도구는 잠재적 취약점을 식별할 수 있으나, 코드의 의미(Semantics)와 문맥(Context)을 완전히 이해하지 못해 다음과 같은 한계가 존재한다.

\begin{itemize}
    \item \textbf{문맥 민감성(Context Sensitivity) 부재}: 동일한 코드 패턴이라도 문맥에 따라 안전할 수도, 취약할 수도 있다.
    \item \textbf{높은 오탐률(False Positives)}: 런타임 제약 조건이나 검증 로직으로 인해 실제로는 안전한 코드임에도 취약점으로 탐지되는 경우가 빈번하다.
    \item \textbf{복잡한 패턴}: 취약점은 단일 구문이 아닌, 여러 코드 요소 간의 복잡한 상호작용으로 인해 발생하는 경우가 많다.
\end{itemize}

Agent 모듈은 레이블링된 데이터셋을 통해 GNN 모델을 학습시킴으로써 이러한 문제를 해결한다. GNN 아키텍처는 코드를 그래프 구조(구문을 위한 AST)로 표현하고, 노드 및 엣지 속성으로부터 특징(Feature)을 자동으로 학습하여 구문적 패턴을 포착하는 데 최적화되어 있다.

\subsubsection{기능 요약}
Agent 모듈은 다음과 같은 기계 학습 파이프라인 전 과정을 제공한다.

\begin{enumerate}
    \item \textbf{데이터 로딩}: 함수별 AST 그래프가 포함된 JSON 파일을 로드한다.
    \item \textbf{그래프 구축}: JSON 데이터를 PyTorch Geometric의 \texttt{Data} 객체(노드 특징, 엣지 인덱스, 엣지 속성 포함)로 변환한다.
    \item \textbf{모델 학습}: 학습 데이터를 사용하여 함수를 취약(1) 또는 안전(0)으로 분류하는 GNN 모델을 학습한다.
    \item \textbf{모델 평가}: 테스트 데이터를 통해 모델의 정확도, 혼동 행렬(Confusion Matrix), 클래스별 통계 등을 평가한다.
    \item \textbf{설정 관리}: 재현성을 보장하기 위해 전체 학습 구성을 저장하고 관리한다.
\end{enumerate}

\subsection{실행 흐름}

\subsubsection{학습 파이프라인}
학습 프로세스는 \texttt{train()} 함수 호출 시 시작되며, 데이터 로딩부터 모델 저장까지의 과정을 수행한다.

\begin{figure}[H]
    \centering
    \begin{lstlisting}[language=Python]
def train(cfg: Optional[TrainConfig] = None) -> None:
    cfg = cfg or _parse_train_args()
    torch.manual_seed(cfg.seed)
    device = torch.device(cfg.device)
    
    # Load datasets and create data loaders
    train_dataset, test_dataset = load_and_split_datasets(cfg.data_path)
    train_dataloader = build_dataloader(train_dataset, cfg, collate_fn=collate)
    test_dataloader = build_dataloader(test_dataset, cfg, collate_fn=collate)
    
    # Initialize model
    dims = infer_dims_from_dataset(train_dataset, kinds=["ast"])
    model = select_model(cfg, dims["ast"][0], dims["ast"][1]).to(device)
    
    # Setup training
    optimizer = torch.optim.Adam(model.parameters(), lr=cfg.lr)
    criterion = nn.CrossEntropyLoss(weight=compute_class_weights(...))
    
    # Train and save
    train_model_from_dataset(cfg, train_dataloader, model, optimizer, criterion, device)
    torch.save(model.state_dict(), weights_path)
    save_training_config(cfg, results_dir, model_info)
    \end{lstlisting}
    \caption{학습 파이프라인 메인 로직}
    \label{fig:train_pipeline}
\end{figure}

주요 단계별 상세 내용은 다음과 같다.

\begin{itemize}
    \item \textbf{구성 파싱}: CLI 인자를 통해 \texttt{TrainConfig} 객체를 생성하여 유연한 실험 설정을 지원한다.
    \item \textbf{데이터셋 로딩}: \texttt{GenericJsonDataset}을 통해 JSON 파일을 PyG 객체로 변환하고, 파일명 키워드 등을 기반으로 레이블을 추론한다.
    \item \textbf{데이터로더 생성}: \texttt{collate} 함수를 사용하여 AST 그래프 배치를 구성한다.
    \item \textbf{차원 추론}: 데이터셋을 검사하여 가변적인 노드/엣지 속성의 차원을 자동으로 계산한다.
    \item \textbf{모델 선택}: AST 그래프를 처리하는 GNN 모델을 인스턴스화한다.
    \item \textbf{학습 및 저장}: 에포크별 학습, 체크포인트 저장, 손실 기록, 최종 모델 및 설정 파일 저장을 수행한다.
\end{itemize}

\subsubsection{평가 파이프라인}
평가 프로세스는 저장된 설정을 로드하여 학습 환경을 복원한 후 수행된다.

\begin{figure}[H]
    \centering
    \begin{lstlisting}[language=Python]
def evaluate() -> None:
    config = load_training_config(results_dir)
    cfg = TrainConfig(**config["training_config"])
    train_dataset, test_dataset = reconstruct_datasets(cfg)
    model = load_model_robust(checkpoint_path, device)
    summary = evaluate_model(model, eval_dataloader, device, cfg.mode)
    save_evaluation_results(summary, "evaluation.json")
    \end{lstlisting}
    \caption{평가 파이프라인 로직}
    \label{fig:eval_pipeline}
\end{figure}

\subsection{컴포넌트 책임}

\subsubsection{GenericJsonDataset}
JSON 파일을 로드하고 PyTorch Geometric 그래프 객체로 변환하는 역할을 담당한다. Pydantic 모델을 사용하여 JSON 스키마를 검증하며, \texttt{juliet\_json\_to\_sample} 변환기를 통해 AST 섹션을 추출하고 그래프를 구축한다.

\subsubsection{Graph Building Functions}
JSON 구조의 노드와 엣지를 PyG 텐서로 변환한다.
\begin{itemize}
    \item \textbf{노드 특징}: JSON의 \texttt{feat} 필드나 노드 자체를 플랫하게 펼쳐 특징 행렬 \texttt{x}를 구성한다.
    \item \textbf{엣지 구성}: \texttt{edges\_ast} 패밀리를 식별하고, 원-핫 인코딩된 패밀리 정보와 엣지 속성을 결합하여 \texttt{edge\_attr} 및 \texttt{edge\_index}를 생성한다.
    \item \textbf{차원 조화}: \texttt{\_harmonize\_graph\_dims} 함수를 통해 배치 내 모든 그래프가 호환 가능한 차원을 갖도록 패딩한다.
\end{itemize}

\subsubsection{Collate Function}
AST 그래프와 가변 메타데이터를 포함하는 샘플들을 하나의 배치로 결합한다. PyTorch Geometric의 \texttt{Batch.from\_data\_list}를 사용하여 효율적인 병렬 처리가 가능한 연결된 그래프(Disjoint Union) 형태를 생성한다.

\subsubsection{Model Selection \& Training}
설정된 모드에 따라 적절한 모델 아키텍처를 선택하고 학습 루프를 제어한다. AST 그래프를 입력으로 받아 취약점 여부를 분류한다.

\subsubsection{Evaluation Functions}
\texttt{evaluate\_model} 및 \texttt{analyze\_sample\_with\_model} 함수를 통해 추론을 수행하고, 정확도, 확신도(Confidence), 오분류 샘플 등의 상세 분석 결과를 생성한다.

\subsection{데이터 변환}

\subsubsection{JSON $\rightarrow$ PyG Graph}
입력 JSON은 AST 노드 및 엣지 정보를 포함하며, 다음과 같은 과정을 거쳐 PyG \texttt{Data} 객체로 변환된다.

\begin{enumerate}
    \item \textbf{노드 특징 추출}: 모든 노드의 고유한 속성 키를 수집하여 특징 행렬을 구성한다 (누락된 속성은 0으로 패딩).
    \item \textbf{엣지 수집 및 특징화}: 엣지 패밀리(타입)를 식별하여 원-핫 벡터를 생성하고, 추가 속성(Guard kind 등)과 결합한다.
    \item \textbf{텐서 생성}: 최종적으로 \texttt{x}, \texttt{edge\_index}, \texttt{edge\_attr} 텐서와 레이블 \texttt{y}, 메타데이터를 포함하는 객체를 생성한다.
\end{enumerate}

\subsubsection{Graph Batching}
배치 처리 시 차원 조화(Dimension Harmonization) 과정을 통해 서로 다른 특징 집합을 가진 그래프들을 정렬하고, 이를 하나의 거대한 그래프 객체(\texttt{Batch})로 병합하여 GPU 연산 효율성을 극대화한다.

\subsection{모델 아키텍처}

\subsubsection{GNN Model Architecture}
본 시스템은 AST 그래프를 처리하는 GNN 모델을 사용한다. GINEStack(Graph Isomorphism Network with Edge features)을 사용하여 그래프를 인코딩하고, 글로벌 풀링 후 선형 분류기를 통해 예측을 수행한다.

\begin{figure}[H]
    \centering
    \begin{lstlisting}[language=Python]
class ASTGNNModel(nn.Module):
    def __init__(self, ast_in, ast_edge_dim, cfg):
        self.encoder = GINEStack(cfg.num_layers, ast_in, cfg.hidden_dim, ast_edge_dim)
        self.pool = global_mean_pool
        self.classifier = nn.Linear(cfg.hidden_dim, cfg.out_classes)

    def forward(self, ast_data):
        ast_repr = self.encoder(ast_data.x, ast_data.edge_index, ast_data.edge_attr)
        return self.classifier(self.pool(ast_repr, ast_data.batch))
    \end{lstlisting}
    \caption{GNN 모델 구조}
    \label{fig:gnn_model}
\end{figure}

\begin{itemize}
    \item \textbf{Encoder}: GINEStack을 사용하여 AST 그래프의 구조적 특성을 학습한다.
    \item \textbf{Pooling}: 글로벌 평균 풀링을 통해 그래프 수준의 표현을 생성한다.
    \item \textbf{Classifier}: 학습된 그래프 표현을 바탕으로 취약점 여부를 분류한다.
\end{itemize}

\subsection{설정}

\texttt{TrainConfig} 클래스는 학습에 필요한 모든 하이퍼파라미터와 데이터 경로를 정의한다.

\begin{itemize}
    \item \textbf{기본 설정}: 에포크 수, 학습률, 배치 크기, 디바이스, GNN 레이어 수, 히든 차원 등.
    \item \textbf{데이터 경로}: JSON 파일 경로 및 레이블 추론 규칙(\texttt{label\_key})을 정의한다.
    \item \textbf{레이블 추론}: 명시적 레이블 필드, 파일명 키워드 매칭("bad", "patched"), 휴리스틱 순으로 레이블을 결정한다.
\end{itemize}

학습 시 \texttt{training\_config.json} 파일이 생성되어 실험의 재현성을 보장하며, 모델 로딩 시 수동 설정 없이도 아키텍처를 복원할 수 있게 한다.

\subsection{분석 파이프라인 통합}

Agent 모듈은 전체 정적 분석 파이프라인의 최종 단계로 통합된다.

\begin{enumerate}
    \item \textbf{CPG 생성}: Joern을 이용해 소스 코드를 CPG로 변환.
    \item \textbf{Template 생성}: CPG를 정규화된 템플릿 노드로 변환.
    \item \textbf{Graph 추출}: 템플릿에서 AST 그래프 추출 및 JSON 익스포트.
    \item \textbf{GNN 분석}: Agent 모듈이 JSON을 입력받아 취약점 여부 판별.
\end{enumerate}

CLI 명령어 \texttt{uv run train} 및 \texttt{uv run evaluate}를 통해 독립적으로 실행 가능하며, 학습된 모델은 배치 예측이나 정적 분석 결과의 필터링 용도로 배포되어 활용된다.